\section{Conclusiones}
\label{sec:conclusions}

Litoteca integra un flujo completo de clasificación de rocas en una aplicación web
móvil: preparación trazable del dataset, entrenamiento en Azure Custom Vision, proxy
de predicción en Cloudflare Pages y una interfaz que presenta la predicción con
contexto geológico. El aporte principal no es solo la etiqueta de clase, sino la
combinación de clasificación, guardrail roca/no-roca y ficha explicativa para un
usuario no experto.

La preparación de datos fue decisiva para la credibilidad del modelo. La
deduplicación por imagen original, la reconstrucción de splits sin fuga y la
exclusión de conflictos de etiqueta redujeron el riesgo de métricas optimistas. Los
resultados verificados de Quick v2 y del detector roca/no-roca muestran un prototipo
funcional, con buen desempeño en varias clases y confusiones concentradas en
familias visualmente cercanas.

Como limitaciones actuales, el prototipo cubre nueve clases ígneas y sedimentarias,
depende de la calidad de la fotografía y del endpoint de predicción, y mantiene
recall bajo en clases difíciles como \emph{shale}, \emph{clay} y \emph{chert}. La
validación sintética por API aporta evidencia reproducible, pero no sustituye una
evaluación de campo más amplia.

El documento se cierra con Quick v3 evaluado: obtuvo P 79.70\%, R 54.58\% y AP
74.64\%, con validación sintética API de 96/146 (65.75\%), sin superar a Quick v2.
Por ello, la entrega se mantiene con Quick v2 como modelo publicado. El
entrenamiento Advanced permanece como trabajo futuro: si completa y mejora el
desempeño verificado, podrá conectarse por debajo del mismo contrato de la
aplicación sin cambiar la experiencia de usuario.
